% =====================================================================
%  R. Nielsen, PERSONAL GROUNDS
%  In: For Idea and Reality.  A Journal, edited by R. Nielsen, Bj. Bjørnson and Rud. Schmidt.  Volume for 1869.
%  Copenhagen: Chr. Steen & Søn; printed by Louis Klein.  pp. 585-622.  (Personlige Grunde, in For Ide og Virkelighed, Aargang 1869.)
%  ENGLISH TRANSLATION, from transcription.tex (Danish).
%  GENERATED by ebook-method/for-idee-1869/tr_build.py --write from this template and the translators' parts in
%  ebook-method/for-idee-1869/tr/ (TRANSLATOR-BRIEF.md; terms: GLOSSARY.md, the titles and four key terms approved by
%  Hans 2026-10-13).  Conventions: TRANSLATION-PLAYBOOK.md.  Page numbers of the 1869 print in the margin.
% =====================================================================
\documentclass[12pt,a4paper,openany]{book}
\usepackage[utf8]{inputenc}
\usepackage[T1]{fontenc}
\usepackage{libertinus}
\usepackage{libertinust1math}
\usepackage[english]{babel}
\usepackage[protrusion=true,expansion=true]{microtype}
\usepackage{setspace}
\usepackage[margin=1.3in]{geometry}
\usepackage{fancyhdr}
\usepackage{hyperref}
\hypersetup{pdftitle={Personal Grounds}, pdfauthor={Rasmus Nielsen}, hidelinks}
\setlength{\headheight}{15pt}
\pagestyle{fancy}
\fancyhf{}
\fancyhead[LE,RO]{\thepage}
\fancyhead[C]{\textit{R. Nielsen: Personal Grounds (1869)}}
\setstretch{1.3}
\usepackage{marginnote}
\renewcommand{\thefootnote}{\arabic{footnote})}   % the print numbers each page's note ¹) (p. 553: *))
\newcommand{\opage}[1]{\marginnote{\footnotesize\textit{[#1]}}}
\newcommand{\essayhead}[2]{\begin{center}{\large #1}\ifx&#2&\else\\[0.3em]{\normalsize #2}\fi\\[0.5em]\rule{3em}{0.4pt}\end{center}\medskip\noindent}
\newcommand{\signature}[1]{\par\smallskip\hfill #1\hspace*{2em}\par}

\begin{document}

% the volume's title page (PDF 7)
\begin{titlepage}
\begin{center}
\vspace*{3em}
{\Huge For Idea and Reality.}\\[1em]
\rule{8em}{0.4pt}\\[1em]
{\large A Journal,}\\[0.8em]
{\small edited}\\[0.4em]
{\footnotesize by}\\[0.6em]
{\large R. Nielsen,}\\[0.4em]
{\large Bj. Bjørnson and Rud. Schmidt.}\\[1em]
\rule{8em}{0.4pt}\\[1em]
{\large Volume for 1869.}\\[1em]
\rule{8em}{0.4pt}\\[4em]
Copenhagen.\\[0.4em]
Chr. Steen \& Søn, Publishers.\\[1.5em]
{\small Printed by Louis Klein.}\\[3em]
{\small Translated from the Danish}
\end{center}
\end{titlepage}

\essayhead{Personal Grounds.}{}
% --- p. 585 ---
\setcounter{footnote}{0}%
\opage{585}Anyone who knows a little about himself and the world
will surely have had the experience that there is a conviction
that is hard to put forward and justify
to others, and this not from a lack of grounds,
no, on the contrary, from an excess of grounds. It holds
in most cases of womanly conviction. A
young girl, for example, may in a matter of the highest importance
be quite clear in her own mind, perfectly
certain of her cause, and yet, when she is to explain herself, it
is as if she sank in a torrent of grounds; she falls silent.
She cannot appeal exclusively to any
single ground, and were she to try to enumerate all
her grounds, she might just as well, she feels,
begin to count the leaves in the forest or the drops
in the sea. That her grounds are personal, not
scientific, is surely a superfluous remark. One
would, however, be greatly mistaken if one imagined
that a conviction from subjective, personal grounds
were suited only to women, not to men. Absolutely
personal grounds hold everywhere where it is a question of
an absolutely personal matter.

In Oehlenschläger's Axel and Valborg it is, as
is well known, Black Friar Knud who, in the circumstance
% --- p. 586 ---
\setcounter{footnote}{0}%
\opage{586}that bride and bridegroom are ``baptismal siblings,'' shrewdly
discovers a ground, fully valid by the law of the Church, for
preventing the marriage; that the ground is objective, a ground
that admits no objections, Axel learns when
he, who ``can also read,'' sees with his own eyes that it
stands in the church register. The objective ground does not convince
Axel that his love for Valborg is sinful
--- in that case the prohibition would have to be conceived from the
religious side, thus subjectively --- no, it does not convince him,
but it shows him a hostile resistance, a barrier,
hard as fate, inflexible as necessity; it
shows him the impossibility of reaching the goal of his longing, his
Eden, ``his blessed Valhalla.'' Thus a single objective
ground can press like the hand of Death on the breast of the living
and subdue all the grounds of the heart, however
heartfelt, however strong these may be.
The tragic element in the catastrophe is an expression of the fact that subjective and objective
grounds are dissimilar.

But such an expression can be had in another
way as well. We may well imagine that Axel's friend Vilhelm,
by one of the many wonderful coincidences that always
stand at the disposal of the poet's imagination, had in good time
got wind of the secret, had come to Rome
a few days after Axel's departure, and had then in all quiet
worked to have the overlooked impediment removed.
Just as Knud has opened the church register, and Axel throws himself
at Valborg's feet with the cry of despair:
``Valborg! Valborg! it is over,'' Vilhelm steps forward,
takes up Valborg's reply: No, no! and goes on: Hadrianus
Episcopus, servus servorum Dei,\dots{} to be brief:
His Holiness has by special grace permitted that Axel
% --- p. 587 ---
\setcounter{footnote}{0}%
\opage{587}Thordsøn may marry Valborg Immer's daughter, notwithstanding that they
are godsiblings, yet on the condition that the said
Axel Thordsøn first show with clear, irrefutable grounds
why he loves Valborg precisely and will never
be able to love any other woman. --- The message is
a joyful surprise to the couple so lately unhappy;
Axel asks in rapture: Is that all,
does the Pope demand nothing else? Vilhelm. No, nothing else.
And yet the matter is curious enough. It was a
Cardinal who interceded for you, a fine mind, a
subtle Scholastic. Now you must know, what you surely
have not noticed, that among the learned there is carried on
a heated dispute about the relation of the will to motives,
about the relation of human freedom to necessity,
some maintaining that true freedom at bottom
is nothing other than a kind of psychic natural necessity,
while others hold that there is a profound
difference between necessity and freedom, that the free
will is determined by grounds only insofar as it determines
itself. The Cardinal now represented to the Pope that it
would greatly contribute to clarifying the matter if some
among the many lovers who yearly besiege
the Holy Father with petitions for dispensation were called upon
to express themselves more closely about their grounds. For lovers
often speak an equivocal language when they speak of grounds.
At one moment they say that they love from necessity, that they cannot
do otherwise; then it is a matter of life and death, and one must
see to helping the unfortunates as best one can.
But then again they say that they love with freedom, that love
is a free matter, that with full resolution they have
chosen the chosen one, picked the picked one, that they cannot
% --- p. 588 ---
\setcounter{footnote}{0}%
\opage{588}do otherwise because they do not \emph{will} otherwise; but when the choice
thus depends on the will and happens with freedom,
the suppliant could just as well choose another, and there
is no pressing ground to grant him dispensation.
By collecting and comparing the many statements
His Holiness could then learn whether the will in this important point
is moved by ``the dark, the egoistic and instinctive,''
which ``hangs together with the unopened unity of
individual natural being,'' or by ``the universal of knowing, the determination
of the universality of essence, which expresses itself in the
being-for-itself in the form of liberated universality, which
belongs to self-consciousness,'' in other words, whether they really love
from necessity or from freedom, or partly from necessity
partly from freedom, or\dots{} Axel. Do
stop with this long string of learned phrases; you
know I am no cleric. I go straight to the point;
I have a ground, a single simple one; my ground
the Pope will not reject, it is my Valborg herself ---
``fir-grown Valborg, rose of all flowers''! Vilhelm.
Yes, my friend, that is a lyrical ground; that ground you
cannot offer His Holiness; nor was it that kind of
grounds of which the Cardinal spoke, for he meant
objective, scientific grounds. Axel. Then you must
help me to mint the ground with the proper stamp; for
I will tell you, it is with my ground as with Draupner,
Odin's ring: it is only one and yet becomes many.
If the Pope does not accept the one ground, then he must
accept the other. I have many grounds, Vilhelm,
when I think it over, many; ``The sea does not have so many
animals in its waves, never did the night see so many stars
in the vault of heaven, not so many birds dwell
% --- p. 589 ---
\setcounter{footnote}{0}%
\opage{589}in the forest, not so many blades of grass has field and
meadow,'' as my heart has grounds why I love Valborg.
--- So the deliberation begins. The wedding is postponed, Axel
ponders, Valborg waits --- she will have long to wait.

So much is surely clear from this, that a subjective personal
ground can never become impersonal, any more than
the reverse. If the dissimilar grounds could in the scientific
sense melt together into one, then the thinking
consciousness would also be able to convert
the one into the other, to transform the one by means of
the other. But how impossible such a transformation is
can be seen from the examples cited. The objective ground,
that godsiblings, by the Church's prohibition, may not marry
one another, can, when it works in all its objectivity,
at most bring about only that the lovers bow beneath
Providence, that in quiet patience they renounce
the happiness of life. It is a power of humility in the woman's
soul that she can bow more inwardly beneath
the inevitable than the man; we hear Valborg
speak for quiet resignation. ``Quietly shall I sit
thus in a little cell, work gold in silk, live out gently
and sorrowfully my time, like the turtledove''\dots{} As she
consoles herself, so she will also console her friend, that he
shall be patient, quiet. ``You shall travel to your
farm, up to your sister, the good Helfred. Time heals
all things in the world, may it heal your wounds. Nature in her noble,
quiet grandeur, a loving sister's company soothe the heart
more than merry company.'' The subjective has hereby
by no means become objective, but the personal ground
has, by collision with fate, been transformed in itself.

It is surely not Valborg's meaning, when she speaks of,
% --- p. 590 ---
\setcounter{footnote}{0}%
\opage{590}that Time heals all wounds, that she will in time forget
Axel, or that he should forget her. What
changes with time is properly only the outward,
not the inward. Outwardly, in the eyes of the world, she must
be parted from the beloved, but she does not lose him;
she knows a place where he always is, a sacred
place in the inner life. In there, where he always is, there
she builds the church --- with old graves for dear
memories --- there she lives, there she keeps vigil and holds
daily divine service: wreathes the pillar, greets the pilgrim,
weeps at his breast, follows him to the altar,
to the altar of sorrow, to the wedding, the wedding that
no Black Knud interrupts. It is only subjective,
all of it subjective, but it is personal. Were
the subjective, on the contrary, transformed into the objective, the
personal into the impersonal: then the matter would have to take
another turn. It would then have to be that Axel suddenly,
through the blow of fate, had become an objective thinker,
had come to see that the original thing in man's
personal being is not heart and will, but
thinking and knowing; he would then have to explain himself to Valborg.
Understand me, Valborg, the hard blow that strikes
us we have deserved ourselves; we have committed a great
sin, not against God in heaven --- that is a subjective
matter --- but against the eternal laws of thinking and being, against
metaphysics. You see, we began in our youthfulness
with our own being, with ``being in the form of
individuality,'' we regarded the contingency of this our restricted
subjectivity as a reliable expression
of ``individual being in general,'' we have unfortunately
taken the concept of existence in ``an arbitrary
% --- p. 591 ---
\setcounter{footnote}{0}%
\opage{591}double meaning'': that is an unclarity in thought, ``a surreptitious assumption,''
which avenges itself now that we awake to consciousness.
Understand me: ``all being in the form of individuality
is existence; there can thus be such an existence
as is not consciousness, not theoretically. But when \emph{consciousness}
becomes the subject of existence, then
existence includes the theory, ``for consciousness exists only
as being theoretical.'' --- You stare at me, you
do not understand me; give it time! You will surely come to
understand me. That a youth loves a girl can be defended
on objective grounds; it is a manifestation of
the universal of nature; it is the life of the species, continuing
itself through the individuals, whose existence insofar must
be said to fall under individual being in general;
but that two people in love should have the right to deify
one another to such a degree that she despairs and will go
into a convent if she cannot have him, precisely him,
although there are other young men of equally lovable
qualities who offer themselves; and that he becomes desperate
and will go to war and absolutely get himself killed, when he cannot
have her, precisely her, the chosen one: you see, Valborg,
such a choosing cannot be defended on objective
grounds, it is to put ``being in the form of individuality''
for ``individual being in general,'' it is
to violate ``the theory.'' Prove now in deed that you have
grasped the objective, that your consciousness has at last
become your existence's subject,'' and that this your subject
really includes the theory: then you will become Hakon's bride
--- to the honor of the objective! Queen of Norway --- to
the good fortune and glory of the universal! Then I am consoled;
you will be too. Farewell!

% --- p. 592 ---
\setcounter{footnote}{0}%
\opage{592}How Valborg will take the explanation, whether she
will take herself or the speaker to be insane, whether
in raging anger she will rise against the faithless one, or
whether giddily she will sink down in faintness, about that
opinions may differ; but this much is certain: she
will not understand the explanation, will not grasp the objective of the theory.
She gave him her heart, her faith; this one thing she understands,
this subjective thing; but the objective? how
should she understand ``the objective of the theory'': she is after all
--- only a woman!

``Love,'' says S. Kierkegaard, ``has after all
its priests in the poets, and now and then one hears a voice
that knows how to keep it in honor; but of faith nothing is
heard, who speaks to the honor of this passion? Philosophy
goes further. Theology sits painted at the
window and courts its favor, offers its beauty for sale
to philosophy. It is said to be difficult to understand
Hegel, that is a work of wonder, but to get beyond
Abraham, that is a small matter. To go beyond Hegel,
that is a work of wonder, but to get beyond Abraham
is the easiest thing of all.'' Whether Kierkegaard is just
right when he calls faith a passion is a
question in itself; but there is one thing in which he is unconditionally
right, namely this, that faith, just as much as
love, is a personal matter. To wish in this
chapter to overtake such an author with school criticism
on a couple of pages and to give it the appearance that one
goes further, goes beyond Kierkegaard, while
one in reality goes back not merely to Hegel,
but even behind Hegel, is not to the advancement of truth.
It is not for nothing that Kierkegaard in all
% --- p. 593 ---
\setcounter{footnote}{0}%
\opage{593}his writings, not merely in the aesthetic and philosophical,
but also in the religious, edifying ones, points to
the erotic and takes examples from natural relations of love.
He wishes hereby to make it quite conspicuous
how sharp the opposition between the personal
and the impersonal is. For what for life and personality
is the most important, the most essential of all: \emph{this}
human being, \emph{this} soul, this heart, must for the consciousness
``that exists only as being theoretical''
necessarily become something wholly inessential, something
vanishing and therefore insignificant. Theory reaches
only \emph{a} soul, never \emph{this} soul, only \emph{a} heart
in general, the concept of the heart, never \emph{this} heart.
What the heart in and for itself, the objective heart, desires
and by virtue of the concept must desire, is the main thing; what
a single heart, this soul, \emph{this} subjective heart
longs for, is indifferent. If a single heart breaks of
sorrow --- a single fly drowns in a cream pot --- what
power lies in that? There are hearts enough in the
concept of the heart, for the concept of the heart, the universal heart, embraces
all possible hearts.

In this spirit Hegel speaks when in the second part
of his Aesthetics he characterizes the chivalry of the Middle Ages
and romantic love. The choice of subjectivity
appears to him as a wilfulness, a ``stubbornness of particularity.'' ``The collisions of love,''
he says, ``especially when they are set fighting against substantial
interests, always retain a touch of contingency and incalculability;
it is subjectivity as such which, with
its claims not valid in and for themselves, sets itself up
against that which for the sake of its own essentiality has
% --- p. 594 ---
\setcounter{footnote}{0}%
\opage{594}a claim to recognition. The individuals in the high tragedy of the ancients,
Agamemnon, Clytemnestra, Orestes, Oedipus,
Creon, etc., have indeed likewise individual ends,
but the substantial, the pathos which in the content of the action
is the moving force, has absolute
justification and is in itself of universal interest\dots{}
When Clytemnestra's crime is not punished in this concrete case,
when the injury Antigone must suffer as
sister is not avenged, then this is a wrong in and for itself.
But'' --- here he aims at chivalry --- ``these
sufferings of love, these shipwrecked hopes,
this infatuation in general, these infinite pains
which a lover feels, this infinite bliss and
blessedness which he imagines, have no interest in themselves
universal, it is something that concerns only
himself. Every human being certainly has a heart for
love and a right to become happy through it, but
when he here, in just this case, under these and these
circumstances, precisely with regard to this girl,
does not reach his goal, then no wrong is thereby done.
For it is not in and for itself necessary --- dass er sich
gerade auf dieses Mädchen kapricionire, and then we should
take an interest in the highest contingency, in the arbitrariness
of subjectivity, which is without all extension
and universality.''

This conception of romantic love, somewhat Hieronymite, as it might perhaps seem, is
objective, strictly scientific. How genuinely scientific
it is we can easily prove. For the solution
of a geometrical problem, e.g., a point must be marked off
or a line drawn. But every time a point is
% --- p. 595 ---
\setcounter{footnote}{0}%
\opage{595}marked off, it becomes \emph{this} point; the line that is drawn becomes
each time a \emph{this} line. He who on behalf of universality
has a mind to tease the mathematician
may with a certain ground raise an objection. You are to mark off
a point, that is necessary; but you mark off
\emph{this} point, that is surely contingent; you want to draw a line,
but why precisely this line? You promise us the universal
and offer us only the particular, you aim at the
necessary and hit only the contingent: that is a
strange method. The mathematician will, however, easily extricate
himself from the embarrassment. He must indeed concede that
every determinate point, wherever he may mark
it off, never becomes the general point, and that the same
must be said of the line; but the contingency that thus
runs along in the choice of the single point and the single
line is so far from weakening necessity
that it is, on the contrary, itself necessary in order to give the universal
existence. Since now the one point in this respect
can be as good as the other, the one
line as satisfactory as the other: so the
arbitrary, insofar contingent, point and the arbitrary,
insofar contingent, line may very well be regarded
as fully valid representatives of point and line in general.

According to the theory it should now, strictly speaking, stand
with the lover and his possible choice as with the mathematician
and his possible points; but that is impossible.
Among all lovers famed from poetry, the one who most
seems to approach ``the determination of the universality of essence''
with the corresponding ``being-for-itself in the form
of liberated universality,'' is beyond comparison Don Juan.
% --- p. 596 ---
\setcounter{footnote}{0}%
\opage{596}And yet even a Don Juan, with his free choice, his
almost unsurveyable wealth of blooming possibilities,
despite the ease with which he chooses and the speed
with which he changes, can impossibly relate himself to
``baronesses and countesses, peasant girls and princesses''
as a mathematician relates himself to points, surfaces
and curved lines. That he, every time he steps from possibility
into reality, to understand, when on the way of existence
he erotically choosing executes, must be content
with one at a time, a single one, a This, is not the difficulty;
for the same can literally be said of
the mathematician, who must be content with this point, this
line. Don Juan will concede to Hegel that every human being
has a heart for love and a right to become
happy through it, and also that a lover is not for life and
death unconditionally to bind himself to a single one as to the
only one in the world, since it is surely not in and for itself
necessary ``dass er sich gerade auf dieses Mädchen
kapricionire''; and yet Don Juan agrees as little with
Hegel as with Hieronymus: he satisfies only the
subjective, not the objective; instead of bringing the
contingent to merge into the necessary, he leaps
over ``the concept'' as lightly as he leaps over the churchyard wall;
he parodies the theory and brings every moment the necessary
to merge into the contingent, and this although
``the arbitrariness of subjectivity'' has in him precisely
what Hegel demands, ``extension and universality.''

The erotic choice, whereby inclination and passion
often gain such a preponderance over cold reason
that the heart runs away with the understanding, is surely
--- one says --- of so subjective, individual a nature that
% --- p. 597 ---
\setcounter{footnote}{0}%
\opage{597}the personal here in many respects collides with
the pure concept, for which reason the universal of the theory, as concerns
this delicate chapter, prefers to confine itself to
being delimiting without being grounding. But
as we come higher up into the moral regions,
the objective will come into its own; the voice of duty
will rise and impress ``the universality of individual
being.'' One need only think, e.g., of the duty
to love one's kind, one's people, one's fatherland. --- Here we meet
again the old error of the theory.

The theory demands that true personality shall
bow to the universal of existence and take its point of departure
in the objective, the necessary. That personality
will not do, it protests; its highest law is not
objective, but subjective: it has its point of departure in itself.
Only with this point of departure, by virtue of this law,
the law of self-determination, of the will, of self-development, of freedom,
does it recognize the objective and satisfy the
universal. The law of gravity is objective, but the law of duty is
subjective, the former is impersonal, the latter personal. It is
duty to give up one's wilfulness and to master the self-loving
inclinations of one's heart; it is duty to love one's
kind, one's people; it is duty to sacrifice oneself for one's fatherland.
But however different the duties may be, the law of duty is
still one and the same, it is an inner law of
personality, an inviolable law of the self. Personality
gives itself up; why? Because in self-surrender it
feels that it does not lose, but gains --- itself! Personality
sacrifices itself; why? Because in the moment of
sacrifice it recognizes that it is just now satisfying its
original drive, its innermost, truest, deepest need.
% --- p. 598 ---
\setcounter{footnote}{0}%
\opage{598}When an Arnold Winkelried in the press of battle seizes
the hostile spears all around, gathers them in his breast and
presses forward: then it happens not from scientifically objective,
no, but from subjective, unconditionally subjective, personal
grounds. That the hero has breathed in the objective and is moved
by the universal, in that he is moved by the Idea, anyone
can understand; but the moving ground is subjective:
it is in the will, for it is in the heart, there where
no theory reaches. Were the theory to make it comprehensible
on objective grounds why precisely this single individual,
this one, in the extremity of danger sacrifices himself for all the others:
it would become a long-winded deliberation. He loves
freedom, he loves the fatherland, he loves those
dear to him at home; but the ground, the objective ground?
It is just as impossible to state the universal ground
why a Winkelried has come to love his free Switzerland
so exclusively, so deeply and inwardly, as it is
impossible to find out the universal ground why Romeo can love
only his Juliet, or Axel his Valborg.

That the ideas belonging to so-called objective ethics
stir just as originally in mind and heart as
the mutual feeling of love of the loving individuals, of that
our cited tragedy affords precisely a telling
example. A way opens to circumvent the objective prohibition;
Axel's friend has prepared everything for flight;
at parting the lovers kneel on Harald Gille's grave:
then notes of a horn are heard far off. Axel knows
these notes; it is ``a warrior's lur.'' To Vilhelm's exhortation:
``follow where your good Norn beckons,'' he answers:
``my good Norn beckons me to duty.'' Now the movement of the idea
begins personally from a new point of departure;
% --- p. 599 ---
\setcounter{footnote}{0}%
\opage{599}but Axel is instantly at one with himself: here there is no
strife at all between objective and subjective; only an
inner transformation, an ethical conversion of personal
grounds. ``Now I cannot travel; it is impossible, flatly
impossible, Valborg! Vilhelm. You will? Axel. Defend
Hakon Broad-shoulders. Vilhelm. Your enemy? Axel.
What enemy? Him I have defiantly turned my back on;
now my King is beset in distress, to him I owe my
life. Vilhelm. You owe your life to the welfare of the fatherland, not to
Hakon. Do you mean to benefit Norway by defending him,
the young fool, who transgresses all duties for his passion?
Who scorns virtue and the subject's right?
Who forgets wholly the whole for his own advantage. No,
Axel! who would be king, it behooves him to think,
act kingly. Come! Erling is a renowned and valiant
chieftain, noble. Do not delay his victory; he will not
misuse it. Let Norway's scepter gleam from the hero's
hand; it rusts in a nithing's.'' --- It is a peculiar
turn the matter has now taken. So long as the hindrance lay in
the canonical prohibition, it might seem as if the conflict
between the personal and the impersonal arose
from this, that inclination strove against duty, as if
it were only subjective passion that wanted the objective
away, because universal, in themselves valid grounds stood hostile
against the heart's dearest wish. Here the relation is now
precisely the reverse; now it is precisely the objective
that is on the side of inclination and favors
the wish. In and for itself, that is to say: in the objective sense, there
are no compelling grounds why Axel should stay and
take part in the battle. Why should he do so, indeed?
It cannot be for Hakon's sake, for in and for itself
% --- p. 600 ---
\setcounter{footnote}{0}%
\opage{600}Hakon has great faults, and Erling great merits. It cannot
be for the fatherland's sake, for if it may be
presupposed that Erling will possibly become a better king
than Hakon, then the fatherland has in and for itself
not lost but gained, if Erling is victorious. But if it is possible,
indeed even probable, that Axel would rather harm than
benefit Norway by staying and fighting for Hakon: then
neither can he be obligated to it for honor's
sake, for to harm one's fatherland in the imagination that
one benefits it is in and for itself no honor. The matter is
at all events so politically questionable that Vilhelm, at least
for the time being, must seem to have the whole right on his
side; with a small change we may apply his words:
``When all has succeeded, when the high Schwarzburg
with its ring wall safely protects us; when Father Erland has in
the castle chapel, where Leonora was forced to Rudolf, laid
your hands in one another, there will be time enough'' --- for
deliberation. Where, then, is the truly decisive ground for Axel?
It is in his own self, his strong sense of honor,
his chivalrous cast of mind; the knight recognizes
his duty: the ground of the obligation is personal.

If we now set the patriotic against the erotic, so that
the former denotes the extreme objective, the latter the extreme
subjective, then both extremes must necessarily proceed from
the moral self; both have their center in the idea of
morality, and the idea of morality in personality.
In sharp opposition to the arbitrariness of individuals, their contingent
desires and whims, the idea of morality certainly has
its inviolable laws of duty, and in the universal validity of the laws of duty
its objective; but it must be impressed in the most emphatic way
that these laws of duty all without
% --- p. 601 ---
\setcounter{footnote}{0}%
\opage{601}exception are laws of the self, thus in the strictest sense
laws of freedom, that the obligation comes from within, from the heart's
ground, from the will; for the will is in the heart, and as the
heart is, so is the will. Kant has rightly
emphasized that respect for the command of duty is inseparable
from respect for oneself, that the moral law in man's
breast must, like the starry heaven above his head, awaken
his highest admiration.

Could it now rightly be said that man's self
is its own principle, that it, as atheistic ethics
presupposes, forms itself exclusively through itself, develops
itself, sustains itself, etc., then it would undeniably
follow that our duties to our neighbor
without more ado merged into duties to ourselves; and the reckoning
would surely show that everyone is nearest to himself. But
it is precisely a fundamental contradiction to want to assert the
human self by excluding the divine.
``The idea of the good is the idea of the will: to the absolute
good corresponds the absolute will, the absolute self.
The absolute self is the holy God. Over against
the absolute self stands then, with the world as intermediary,
man, the relative self, the self that
is not good, but obligated to the good, the relative
will, which is not holy, but sinful''\footnote[1]{On Obstacles to and Conditions for Spiritual Life. P. 101.}. Here
is the turning point where morality points back to
religion as to its original foundation: the principle of morality
stands and falls with religion. That religion
must bear morality can be pointed out, but
not proved; here is the fact which
% --- p. 602 ---
\setcounter{footnote}{0}%
\opage{602}science can illuminate, but not ground; and why not?
Because the principle is personal.

Were an objective demonstration with corresponding
grounding to be possible, then the principle of contradiction,
well known from logic, would have to provide the guideline.
But anyone who knows a little of cultural history
must be able to tell himself that matters by no means stand
with views of life as with scientific theories:
the concepts of spirit always bend easily to the view of life
once chosen, but never the reverse. Voices are heard
crying: ``We have no place for a willing
God, we feel no need of personal immortality, we
pay no heed to subjective grounds, we do not like the inverted
science,'' etc. That is to say, in other words,
that there are thinking individuals whose ethical self-feeling
is so extremely weak that in the midst of the objective world
they cannot perceive the independence of subjectivity, but
regard the laws of freedom in the spirit as a theoretical reflection
of the laws of necessity in nature. Against this
there is from the standpoint of knowledge nothing to object. He who
absolutely will derive his self from the selfless and bow
personally beneath the impersonal, let him bow,
let him derive, that is his own affair. To reject
religion everyone in knowledge should be free to do; but
when someone in his self-deception goes so far as to want, in the
name of science, to make the world believe that he really
can refute the principle of religion, then protest
must be made; for to reject is not the same as to
refute.

The harshest opposition and yet most immediate
confusion of subjective and objective, personal and
% --- p. 603 ---
\setcounter{footnote}{0}%
\opage{603}impersonal grounds we find in the controversies about the historically
existing religions. A Muhammadan, a Jew,
a Christian would readily agree that there can be only
one single true, saving faith, that the truth of religion
is the universal truth, valid in and for itself, that
all human beings without respect of persons must unconditionally
bow before the one truth, revealed by God himself and therefore
undeceiving: the decisive weight falls
insofar on the universally valid, the objective. But
when it is then asked which revelation, which faith,
is then in and for itself the unconditionally true, then
agreement dissolves into disagreement. Each of the three has his
religion and his faith; the one is just as certain
of having the truth on his side as the other. The
zealous professors of religions may contend with one another;
they may in seething passion murder one another,
but convince one another on objective grounds they
cannot.

If it is forgotten in the heat of the struggle that the personal
truth of faith must begin with the self, the subjective, and
be appropriated with freedom, then passion at once thrusts itself
in under the objective, the impersonal, then
the one will compel the other by objective proofs,
by objective power, and then we have \emph{fanaticism}\footnote[1]{``The fanaticism of faith is known well enough; its mark is that the self-seeking subjectivity in passionate blindness will transform the subjective truth into an objective power, a law of existence, before which all other subjectivities shall unconditionally bow. The fanaticism of knowledge is less known; but not therefore less characteristic. It naturally has not its origin in science, but in a selfish distortion of science. Instead of taking the cognition of the personal for what it is, a subjective practical cognition, the knower makes self-knowledge without more ado objective theoretical. That is a great psychological self-deception. The knower imagines that the subjectivity from which he proceeds is an objective universally valid being, an undeceiving metaphysics-consciousness, a true normal subjectivity, whereas in reality it is only an abstraction of his own individual subjectivity's contingent state. Far from being freed from `egoism,' the knower becomes by this self-deception on the contrary beguiled by egoism; for to make one's own subjectivity's contingent inclination the objective standard for all other subjectivities is egoism. What normal subjectivity, `by virtue of the objective principle of existence,' has need of, that every human subjectivity shall now essentially have need of, namely to become metaphysical; what normal subjectivity feels no need of, a personal God, a comfort in prayer, an eternal life, etc., such dreams and subjective old-wives' fancies no true --- objectively true --- human subjectivity may in future feel need of. To feel need of such things is to deviate from the normal, to deviate from the normal is to deviate from the universal, to deviate from the universal is to deviate from the truth: the standard is objective. With his objective standard the normal subjectivity now needs merely --- in all impersonality to get a little attack of personal zeal for truth, and the smouldering fanaticism of knowledge will betray itself in attacks on faith.'' --- Mr. Professor Brøchner's Philosophical Criticism Reviewed by R. N. P. 18---19.}.

% --- p. 604 ---
\setcounter{footnote}{0}%
\opage{604}If, on the other hand, the individual, the subjective is held fast
with so one-sided a freedom that the universally valid
content of personal truth falls away, then the difference between
the positive religions is inessential, and then we have indifferentism.

The scholastic orthodoxy of the sixteenth and seventeenth centuries gave Protestant
fanaticism of faith a
theological foundation; to undermine this foundation was
a task for rationalism, but the solution of the task
led to indifferentism. The most spirited representative
of rationalism in the last century is beyond comparison
G. E. Lessing.

The work in which Lessing has combated fanaticism and the orthodox
% --- p. 605 ---
\setcounter{footnote}{0}%
\opage{605}narrow-mindedness with the greatest effect is surely his ``\emph{Nathan der Weise}.'' The fable
is set in the time of the Crusades. Saladin has conquered
Jerusalem; a young Templar, who is in love with Nathan's
adopted daughter Recha, turns to the Patriarch
for counsel. As an embodiment of the Church's
authority, exalted above all reason, the Patriarch demands
that his counsel be followed in blind faith, and therefore comes
into some embarrassment, since he cannot get
to know whether the case before him: that a Jew has with
fatherly tenderness brought up a Christian child in love
of all that is good, is a fact or merely ``a problem.'' The Templar,
who only wished to hear the reverend
man's opinion in a general way, thinks that it might be the
same thing. ``The same thing''? exclaims the Patriarch; ``there
you see, sir, how the proud human reason can
go astray in spiritual matters! --- By no means!'' The Patriarch
has neither time nor inclination to concern himself
with problems of thought --- the pro et
contra of investigations of reason he refers to the theater --- he scents a palpable
fact and in holy zeal wants at once to lay
hold of the Jew. ``Papal and imperial law determines,
does it not, that a Jew who seduces a Christian to apostasy
shall go to the stake. All the more then a Jew who has
done violence to a Christian child.'' But, remarks the Templar,
when the Jew has now out of compassion taken
the child in, which otherwise must have perished in
misery? --- ``Say nothing! The Jew shall be burned. Better
that it had perished in misery than that it was thus saved
to its eternal ruin. Besides,
what right has the Jew to anticipate God? God can,
when he wills, save without him.'' --- Also in spite of him,
% --- p. 606 ---
\setcounter{footnote}{0}%
\opage{606}I should think, make blessed. --- ``Say nothing! The Jew
shall be burned.'' --- It goes to my heart, continues
the young man, especially as it is said that the girl was not
brought up in any definite faith, but has only heard of
God what reason teaches. --- ``Say nothing! The Jew shall
be burned.'' Since the Templar will not make a confession
and say outright who the person concerned is, the
Patriarch wants to go straight to the Sultan; for Saladin has, according to
the capitulation, sworn to protect the Catholic religion
with all its holy rights.

In a few but lively strokes Lessing has here sketched
a picture of a genuine fanatic; whether the fanatic happens
to be Patriarch in Jerusalem or perhaps ``Hauptpastor
in Hamburg'' makes no difference to the matter, so long as the drawing
is correct. --- Blind authority, blind faith,
and blind zeal are so intimately connected that from
any two of these marks one can always infer the
third. When the Church with its infallible dogmas is set up
as the object of unconditional submission in faith, the
subjective dwindles away immediately into the objective,
the personal into the impersonal. Whoever does not believe
in the infallibility of the Church is eternally damned. Now
the dread of an eternal damnation can surely be said to be
a subjective ground; but what is it that awakens this
dread? The authoritative word of the Church, a word more objective
than the sun in the heavens. So self-development comes to a halt at
the first step; for the first step of faith is that a human being
give the Church his eternal soul, his innermost self, his conscience.
Every closer explanation, every ground, every
determining Why is to be sought in the doctrine of the Church, its
promises, commandments and prescriptions: the believer retains
% --- p. 607 ---
\setcounter{footnote}{0}%
\opage{607}nothing for himself except blind devotion, unconditional
obedience. But precisely through this self-surrender,
or rather self-abandonment, religious subjectivity becomes
an organ, a will-less and yet strong-willed instrument
of the divine in itself, the objective. The believer's
cause is now without more ado made God's cause, his will
God's will, his friends God's friends, his
enemies God's enemies. Then passion is kindled, and
the raving of faith begins; for what punishment can well be
too hard when it concerns God's enemies? They must be
persecuted with fire and sword, flayed with whips, crucified,
burned: every means for their extermination is pleasing
to God. The Patriarch in Jerusalem would have nothing
against it if Sultan Saladin, though he himself is an unbeliever
who has hundreds of times deserved the stake, would lend
the holy Church his strong arm and be of assistance to it
in burning all Jews.

Religious fanaticism is an unbridled egoism of faith.
In the same degree as self-understanding is lacking and the personal
grounds are lost in blind submission to a
pretended, external divine authority, the
selfish passion rises. The natural self catches fire.
The phenomenon of fire is a sign of a psychic combustion in which
the impersonal and the personal mutually inflame and consume
one another. The excited self-will imagines that it does
God's will and fights a personal fight for the personal
God. A terrible self-deception! The fanatic's God,
the objective and yet groundless ruler, is in his blind
arbitrariness a selfless self, a pitiless power that
annihilates freedom, a poisonous spirit that infects personality.
This fanatical nature, this principle hostile to
% --- p. 608 ---
\setcounter{footnote}{0}%
\opage{608}civilization and humanity, equally dangerous and to the
nobler self-esteem equally revolting, whether it, as
in the time of the Crusades, breaks out in wild, consuming flames,
or under the mask of Protestant state-church
orthodoxy is blown up as a smouldering ember by dogmatic
inquisitions and vents itself in civil vexations,
this monster of dogmatism, against which Lessing
surely, skillfully in many a contest had swung the blade
of criticism, is what he with aesthetic weapons wants
to get at in his \emph{Nathan} the \emph{Wise}.

The conversation between Nathan and Saladin concerning
the relation among the three religions, the Muhammadan,
the Jewish and the Christian, is highly ingenious; we
think of the parable of the three rings. Nathan relates.
In days of old, in the gray antiquity, there lived
a man in the East who from a beloved hand possessed a
ring of invaluable worth. The stone was an opal that
played in a hundred beautiful colors, and had the secret
power of making everyone who wore it with full
confidence pleasing to God and men. No wonder
that this man in the East never let it come off his
finger, and made the arrangement that it should be preserved
in the family forever. Namely thus: He destined
the ring for that one of his sons whom he loved
most; and determined that he in turn should leave it
to that one of his sons who, without regard to birth, was
dearest to him. The son who had the ring was
the head of the family. So this ring came from son
to son, at last to a father who had three sons,
who were all three equally obedient to him and consequently equally dear
to his heart. It happened then that from time to time, as
% --- p. 609 ---
\setcounter{footnote}{0}%
\opage{609}he talked alone with each of the sons severally,
and each time found that the one with whom he talked confidentially
stood nearest his father's heart, he came to promise the ring
to all three. When he now on his deathbed had
only one ring, but three sons, and could not bring himself to
disappoint any of his dear ones, he sent secret word
to an artist and had him make two other rings, which
in every particular perfectly resembled the genuine one.

Saladin, who has listened somewhat impatiently to this
``fairy tale,'' is taken aback when Nathan declares to him that the right
faith is just as unrecognizable as the genuine ring. --- What,
says the Sultan, is that to be an answer to my question?
Do not make a fool of me; I should have thought that the
religions I have named to you are surely distinguishable,
that the difference even extends to dress, food
and drink. --- But, replies Nathan, as regards the grounds
there is no difference. For do they not all
ground themselves on history, oral or written? But history
must surely be accepted on faith and trust? Whose word and
faith does one then least call in doubt? Surely one's own people?
surely theirs, whose blood we are? theirs, who from our
childhood on have given us proofs of their love?
who have never deceived us, except in cases in which it
was serviceable to us to be deceived? How can I believe
my fathers less than you yours? Or conversely. Can
I demand of you that you declare your forefathers
liars, in order not to contradict mine? Or conversely.
The same holds of the Christians. Is it not so? ---
This argumentation makes a strong impression on Saladin;
he exclaims: By the living God, the man
is right. I must be silent.

% --- p. 610 ---
\setcounter{footnote}{0}%
\opage{610}The question of the objective foundation of religion
thus coincides with the question of the credibility of historical
facts and then particularly with the relation
between religious and historical faith. In a little treatise:
\emph{Über den Beweis des Geistes und der}
\emph{Kraft}, Lessing had already set forth the proposition so characteristic of his
point of view, that contingent
truths of history can never become proof of necessary
truths of reason. It is now remarkable to see
the confusion of impersonal and personal grounds.

Historical events are externally objective
and belong to the past; their credibility can,
especially when it is a matter of phenomena of the supernatural,
of signs and wondrous works, easily become doubtful.
The highest to which one can attain by this road
is only a certain degree of probability. But mere
probability never gives full certainty; to this is added
that the miracle, what deviates from the order of nature, always
must show itself as the improbable. The judgment on
the objective becomes in this way dependent on feeling,
on mood, thus on subjectivity. And when then
the subjectivity, uncertain in itself, reaches out for the objective,
what is more natural than that everyone, as Nathan
says, holds to what is handed down from the fathers and preferably
believes his own people? What is historically objective is that one is
born, if not just into, then yet toward a definite religion:
the born Muhammadan believes in Muhammad, the
born Jew in the law of Jehovah, the born Christian in Christ.

On these terms individuals of different
confession of faith have nothing to say to one another. If it
is not possible to prove from historically objective grounds
% --- p. 611 ---
\setcounter{footnote}{0}%
\opage{611}which religion is in and for itself the true one, then
the one may, so far, be just as good as the other.
The objective has given way and bent beneath the
subjective. Religion has become a matter of piety,
a matter of the heart, and the choice is free. Instead of laying stress
on the dogmas, one now lays stress exclusively on
the pious disposition, on virtue and good works.
``Well''! so speaks the judge in Nathan's tale to
the three brothers who quarrel over the genuine ring --- ``so let
each strive to bring to light his incorruptible, unprejudiced
love, that the stone in his ring may show
its power; let him come to the aid of this power with gentleness,
with heartfelt forbearance, with beneficence,
with inward devotion to God: I invite
you then to appear here again when thousands of years have
passed; a man who is wiser than I will then sit
on this chair and pronounce judgment.'' Saladin, continues
Nathan, if you feel yourself to be this wise man,
the promised one\dots{} This address makes so mighty an
impression on ``the Ruler of the Believers'' that he rushes
up to the Jew, seizes his hand and exclaims, moved:
``I dust? I nothing? O God''! --- With such
a self-confession the law of forbearance is signed;
the pyre of fanaticism is extinguished: humanity has triumphed.

But with this turn of the matter we have come to
an opposite extreme: in compensation for fanaticism
we have now got indifferentism and with this the sickness of doubt,
skepticism. If it is once generally conceded
that none of the positive religions, which give themselves out
as revealed, is in a position to furnish a historically irrefutable
proof of the authenticity of its revelation: then it follows
% --- p. 612 ---
\setcounter{footnote}{0}%
\opage{612}quite rightly from this that professors of the different communions of faith
must make concessions to one another. A professor,
for example of the Christian religion, must then say to himself:
The ground why I am a Christian is, if I am to
be sincere, properly this, that I was born in a Christian
society of Christian parents; had I been born in a
Jewish society, I would naturally have become a Jew, and
if I had been born in the middle of Turkey, why, then I would have
become a Turk. From such a confession there follows not merely
leniency toward those who believe otherwise, thus tolerance,
but also indifference toward the mutual difference
of the positive religions.

The result of the last century's spiritual
struggles was a decisive break with blind authority,
an emancipation from what is externally handed down in history.
One should now think that such an emancipation
from the external must have led to a deep, penetrating
cognition of the internal, the personal;
but that did not happen. Something else happened. Just
as subjectivity, fluttering away from its old dogmatic cage,
began to soar high above all fixed
edifices of doctrine and, in the enjoyment of the fresh air with
the free prospect, gave itself over to all manner of religious notions,
relying on this, that everyone is saved by his faith,
reason asserted its divine right and
set the objective once more in force, if under a
quite peculiar form. We shall not dwell here on
a detailed characterization of the so-called religion of reason,
but merely single out one point to show
how rationalism, near to hitting the right thing,
missed the truth, in that aiming at freedom
% --- p. 613 ---
\setcounter{footnote}{0}%
\opage{613}it hit only necessity, and steering toward the personal
ran aground on the impersonal.\footnote[1]{Cf.\ Philosophy of Religion pp. 28—35; pp. 88---97.} The basic point at which the objective meets
the subjective is the self. The human self bears in
itself a double principle, a double predisposition, a predisposition for
knowledge and a predisposition for faith. When then the awakening
reason sets self-thinking in motion, it
depends on whether the gaze is originally directed toward
the objective or toward the subjective.

If the gaze is directed toward the objective, then not
only the things in the world, but also the interior of the soul, the
essence of the self, become an object of knowledge. No one who speaks
naturally will say: I believe that I exist, that I feel,
see, hear, think, etc.; he will say with immediate
assurance: I know it. The same holds
when it is a matter of the consciousness of others. No one will say:
I believe that these people who here surround me
are really living, that they too feel, see, hear, think
just as well as I; but he will, provided that he
has not a metaphysical screw loose, say straight out: I
know it. There is a system of general determinations
that belong just as necessarily to the concept of self-consciousness
as sides and angles belong to the triangle.

If the development of concepts is now laid out so that the question
is not at all about your feelings and mine, your
thoughts and mine, your will and mine, your self and mine, but
about human feelings and thoughts in general, about
the universal will, the universal self, then the psychological
concepts assume an objective, scientific stamp,
% --- p. 614 ---
\setcounter{footnote}{0}%
\opage{614}so the light of reason shines in consciousness. In this
illumination the laws of morality too would present themselves
as pure laws of necessity; insofar as there is here talk
of faith in God as moral lawgiver, this faith of reason will
be of so universally valid a character that
it looks to the life as though its content could with almost
mathematical exactness be grounded and demonstrated by the
understanding.

Rationalism has quite rightly seen that what
in all religion and morality essentially matters
is a clear cognition of the human
self, its origin from God, its eternal destination,
its infinite perfectibility; rationalism has
seen that true faith must be something other than
simple acceptance of the incomprehensible on blind authority;
it has seen that the truths of religion ought, like
other truths, to be tested by the principle of self-thinking
and confirmed by insight, by conviction.
But the fundamental error is conspicuous.

The mistake of rationalism is that it takes the shadow
for reality, the self's thought-being for the living
self in flesh and blood. Through this mistake all self-thinking
becomes scientific thinking, and the laws
of thought coincide with the laws of the natural world.
However much the old rationalists strove to
hold fast to God, to virtue and immortality, it
did not help. God, virtue and immortality are categories of personality,
which point back to the independence
of subjectivity; but the independence of subjectivity
that rationalism reached for was washed away by
the current of objective knowledge: only necessity itself,
% --- p. 615 ---
\setcounter{footnote}{0}%
\opage{615}the form of the eternal, the cold universal, remained. It
is, with regard to the doctrine of spirit, no accident, but
scientific consistency, that rationalism dissolved
into speculative philosophy, and speculative philosophy into
atheism.

If atheism is to be overcome, it must happen through the self;
but then the self must be conceived, not as a general
abstract being, but as a living existence,
not as a concept of knowledge, but, as regards personality,
as a concept of faith.

To conceive the self as a concept of knowledge is, we repeat
it, to look away from one's own existing consciousness,
one's own feelings, the thoughts of one's own heart, and to grasp
the essence of consciousness in its pure generality, thus
impersonally. From the essence of consciousness in its generality
it is impossible for thought to reach your and
my real consciousness, impossible to reach from the impersonal
to the personal.

It is otherwise when the human being is conceived
from the standpoint of faith. Here everyone begins with himself.
The task is to cognize one's original destination,
to discover the eternal in one's own self and to assure oneself
of the eternity of the self: the task is to begin personally.
That this is a task of faith can easily be seen.

Is the soul of man mortal or immortal? To
get this question answered, we must first of all make clear
to ourselves where it belongs, of what character
it is, and what grounds they are on whose validity
the answer depends. The question is personal; that
is to say: here one cannot begin with the essence of the soul in
general; such a beginning is false, for it
% --- p. 616 ---
\setcounter{footnote}{0}%
\opage{616}is impersonal. From a false beginning a true conclusion is never
derived; an impersonal premise can not possibly
yield a personal consequence. In order to begin
personally, everyone must begin with his own consciousness
and ask himself: is my soul immortal? Three
answers are possible; the questioner can answer himself
\emph{affirmatively}, he can answer \emph{negatively} and he can
comport himself \emph{doubtingly}. That the three answers, if they are
to have personal significance, must proceed from three individuals,
three different human beings, needs no proof;
for he who in one moment assumes that the soul
is mortal, in the next that it is immortal, is surely uncertain
of himself, thus doubting. It is, moreover,
clear that while no necessary transition can be
made from the universal to the individual, a free, personal transition
can on the other hand be made from the individual
to the universal. For this stands firm: as
the single, real human being conceives his
original self, his own human essence, so
he must necessarily conceive the essence of the human race;
he does not make his contingent individuality the standard
for other individuals, it is the essence of the self that he makes the
standard; but the essence of the self he cognizes in his own
self: the rays of light from the universal are refracted in the prism
of individuality, gather in the focal point of personality.

If we turn to the two individuals with the two decisive
answers and test to what extent the one by objective
grounds is able to convince the other, we come
at once to a negative result. He who feels himself absolutely
assured of the immortality of the soul, in other words,
of the eternity of the self, has the ground of this assurance
% --- p. 617 ---
\setcounter{footnote}{0}%
\opage{617}in himself and can not possibly fetch it from elsewhere.
What has influenced him from without, what dispensations
he has experienced, what he has lived through in himself,
what struggle he has fought, what crises he has
passed through before he attained the assurance, of that
the question is not; the question is of the peculiar kind
and nature of the assurance attained. When this
standpoint is held fast, then it is indisputable that the certainty
is self-certainty, a certainty about the self, which only the self
gives, and all corresponding grounds personal.

But he who, in contrast to this, holds himself assured
that the soul is mortal, that it is formed with
the body, develops with the body, dissolves and perishes
with the body, that the human self, even if it
is bearer and point of passage for eternal ideas, is not
itself eternal, should he perhaps not be able to build on
science and justify his view of life by objective
grounds? Impossible! All his proofs, however acutely
he may separate them, and however profoundly he
may gather them, unite into one single great proof, the proof
that he finds himself in a great psychological
self-deception. The deception does not show itself directly in this, that
he assumes the soul to be mortal; no, psychologically seen,
the original self-deception consists in this, that the deceived
imagines himself able to justify his view
scientifically.

He who is to judge of the laws of nature must look away
from his own self, in order to judge scientifically;
he who judges of the triangle must look at the triangle,
not at the self, for the triangle is surely not the self; he
who judges of the elements must look at the
% --- p. 618 ---
\setcounter{footnote}{0}%
\opage{618}elements, not at the self, for the self is no element;
likewise when judgment is passed on the development of the earth, on the phenomena
of the primeval world, etc.: all such judgments are
objective in that the judging person looks away from himself
in order to look at the object.

But when the object of the judgment is the judging person's
own self, then the standpoint must be changed.

The judgment of an object must always conform to
the nature of the object and every irrelevant
consideration be kept away. To investigate an element
such as sulphur the chemist must indeed bring
it into manifold interaction with other substances
and regard it in many combinations, but through all this
it is still the peculiarity of sulphur, its physical
and chemical properties, that he has uninterruptedly before his eyes.
If we now put the self in place of sulphur, then it is surely
evident that he who will judge of the self
must uninterruptedly have his eye fixed, not on the pure
metaphysical self --- for when existence is in question,
the metaphysical self is a phantom, --- but on the real
self, its powers and essence, its spiritual demands.
But the only real self the observer can really
investigate is, without contradiction, his own. The real
self that is the object of the judgment is also the judging
subject.

The self judges itself either to life or to
death! thus the psychological deception is lifted;
the judgment is subjective: the grounds are personal.

The psychological error of rationalism and of speculative philosophy
now lies clear to the day. One wanted,
quite rightly, to begin with self-thinking and go
% --- p. 619 ---
\setcounter{footnote}{0}%
\opage{619}back to the self; but confusing the general concept of the self
with the real self, one grasped knowledge and meant
faith, and the consequence was that the immortality of the soul dwindled
away into the eternity of reason, and the right of personality was
lost in the impersonal.

Those who acknowledge, and those who deny, the immortality
of the soul can not convince one another psychologically,
since their self-cognition is so fundamentally different
that they, with regard to the question of immortality,
must be said to have each his own psychology; but the opposition
becomes far sharper still when the talk is
of religion and morality.

In order to know oneself in the moral sense, a
human being must will himself; in order to will himself, he must
choose himself. The judgment on the reality of the self is
in the strictest sense a judgment of will. He who acknowledges
the immortality of his soul has chosen himself in the eternal
sense, has by this choice not merely found the eternal
in the self, but even, literally speaking, discovered the eternity
of the self. Conversely, he who declares the soul
mortal, provided that he does so with full assurance,
has attained this assurance by an inner, original self-determination, an act of will, a choice. Then only
the doubter remains; but the doubter, who gropes between
objective and subjective grounds, is no personality;
he is a loose, enfeebled subjectivity that runs
astray in itself.

If one will now only grant us that life in
itself is not merely something different, but absolutely different
from death, then one will surely also be compelled to
grant that the judgment that makes the self immortal, and
% --- p. 620 ---
\setcounter{footnote}{0}%
\opage{620}the judgment that makes it mortal, are absolutely different
judgments, and that the difference rests on this, that the
judging subjectivities have submitted to absolutely
dissimilar principles.

All cognitions that concern knowledge and science
belong without exception under the principle that transposes
the existing self into the general essence of thinking,
and which we therefore call the \emph{principle of reason}; all cognitions
that concern religion and morality belong, on the contrary,
under the principle that puts the universal into the
existing self and makes the living human being's own
soul the unswerving center, and which we therefore, when
eternal life is in question, call the \emph{principle of faith}.

That these two principles are not merely dissimilar, but
absolutely dissimilar, can be seen from the corresponding grounds;
for all grounds of reason are in their general validity
objective, impersonal, and hold by virtue of necessity;
but the grounds of faith are in their individual validity
subjective, personal and correspond to freedom. He
who still will hold that there is only a relative, not an
absolute difference between the two principles, let him try
to derive freedom from necessity, the personal from the
impersonal, and he shall surely founder in the attempt.\footnote[1]{If one is to patch together a personal view of life for oneself by borrowing a little from faith and a little from knowledge, a little from the objective and a little from the subjective, a little from the impersonal and a little from the personal: then the view of life will be accordingly.}

Faith in the eternity of the self presupposes faith in
the existence of God; for if the human self did not
have its origin from the divine; if
eternity, omnipotence, omniscience, holiness, blessedness did not
% --- p. 621 ---
\setcounter{footnote}{0}%
\opage{621}belong to the being who is the author of all that exists; if
the human soul were only an organic natural product:
how then should the soul withstand decay,
and the self outlive death? But the conviction
of the existence of a personal God is in the deepest
sense personal.

Just as surely as the cause of personality, so long as there is
a remnant of seriousness left in the race, can never become an indifferent
cause, just so surely must the principle of personality be said
to contain an antidote to religious indifferentism; for
everyone must work for his own, and the reckoning concerns
the conscience. But however many clashes may
arise between individuals with different self-cognition
at different stages of development, in that each
in his own way is zealous for the solution of personal tasks, still
it will for the future not be possible to forget that the law of
personality is the law of freedom, its life a life of
freedom, and the principle itself will in this respect
contain an antidote to religious fanaticism.

True faith and true personality are inseparable.
All that concerns faith concerns personality;
what does not belong to the grounding of the personal
life has nothing to do with faith. By holding
fast to this principle professors of the
different religions will work their way forward to one and the same
confession of faith; for the personal stages of life are different,
but the eternal life of personality is in essence the
same.

Exalted above all soul-darkening strife there rises
in the light of faith the Revealer of truth, the Transfigured One of love,
the eternal prototype of personality. To know
% --- p. 622 ---
\setcounter{footnote}{0}%
\opage{622}the ground of the self by knowing this prototype, to believe
in the eternity of the self by believing in this prototype is,
according to the parable in Nathan, to have the genuine ring. When
then the Judge comes --- even if he tarry a thousand
years --- then he will surely not judge by external signs,
by blind authority according to objective dogmas; but he
will ``reveal the counsels of the hearts'' and judge in such a way that
everyone under the judgment must judge himself and give
him right: he will judge personally, by virtue of personal
grounds.

\signature{R. \emph{Nielsen}.}

\end{document}
